% BHU PYQ -- Philosophy: Introduction to Logic (2025-26 NEP)
\documentclass[11pt,a4paper]{article}
\usepackage[a4paper,top=2.5cm,bottom=2.5cm,left=2.8cm,right=2.8cm,headheight=14pt]{geometry}
\usepackage{amsmath,amssymb}
\usepackage{fontspec}
\setmainfont{Kohinoor Devanagari}
\usepackage{microtype,fancyhdr,enumitem,hyperref,lastpage,parskip}

\hypersetup{colorlinks=true,urlcolor=black,linkcolor=black,pdftitle={PHIMJ32: Introduction to Logic -- BHU NEP 2025-26}}

\pagestyle{fancy}\fancyhf{}
\renewcommand{\headrulewidth}{0.4pt}\renewcommand{\footrulewidth}{0.4pt}
\fancyhead[L]{\small \textit{Banaras Hindu University}}
\fancyhead[R]{\small \textit{PHIMJ32: Introduction to Logic (2025-26)}}
\fancyfoot[C]{\small Page \thepage\ of \pageref{LastPage}}

\newlist{parts}{enumerate}{2}
\setlist[parts,1]{label=(\alph*),leftmargin=*,itemsep=4pt,topsep=3pt}
\setlist[parts,2]{label=(\roman*),leftmargin=*,itemsep=2pt,topsep=2pt}

\newcommand{\pts}[1]{\hfill \textit{[#1]}}

\begin{document}

\begin{center}
  {\large\bfseries BANARAS HINDU UNIVERSITY}\\[2pt]
  {\normalsize B. A. (Hons.) Semester III Examination, 2025-26 (NEP)}\\[6pt]
  \rule{0.8\linewidth}{0.5pt}\\[6pt]
  {\large\bfseries PHILOSOPHY}\\[4pt]
  {\large\bfseries Paper Code: PHIMJ32 : Introduction to Logic}\\[4pt]
  {\normalsize Time: 2:30 Hours \hfill Full Marks: 70}\\[2pt]
  {\small REF NO. 2025-26/D-III/091}
\end{center}
\rule{\linewidth}{0.4pt}
\smallskip

\noindent \textit{Instructions: Attempt all three Sections A, B and C according to instructions given. Write your Roll No.\ at the top immediately on receipt of this question paper.}

\smallskip
\noindent \textit{दिये गये निर्देशों के अनुसार खण्ड अ, ब तथा स के उत्तर दीजिये।}

\bigskip

\section*{SECTION -- A / खण्ड -- अ}
\noindent \textit{Note: Long Answer Type Questions. Attempt any \textbf{three} questions, each in 300 words. Each question carries 10 marks.}\pts{3 $\times$ 10 = 30}

\noindent \textit{किन्हीं तीन प्रश्नों के उत्तर दीजिये, प्रत्येक उत्तर 300 शब्दों में हो। प्रत्येक प्रश्न 10 अंकों का है।}

\begin{enumerate}[label=\textbf{\arabic*.},leftmargin=*,itemsep=12pt]

  \item Explain deductive and inductive arguments with suitable examples.\pts{10}
  
  निगमनात्मक तथा आगमनात्मक युक्तियों की उपयुक्त उदाहरणों के साथ व्याख्या कीजिये।

  \item Discuss all kinds of truth-functional statements using truth table method and suitable examples.\pts{10}
  
  सत्यता सारणी विधि एवं उपयुक्त उदाहरणों का उपयोग करते हुए सभी प्रकार के सत्यता-फलानात्मक कथनों की व्याख्या कीजिये।

  \item Determine the validity of the following arguments using the Truth Table Method:\pts{10}
  
  सत्यता सारणी विधि का उपयोग करते हुए निम्नलिखित युक्तियों की वैधता निर्धारित कीजिये:
  \begin{parts}
    \item 1.\ $A \supset B$ \quad 2.\ $(A \lor C) \supset D$ \quad $\therefore A \supset D$
    \item 1.\ $(M \lor N) \supset (D \cdot E)$ \quad 2.\ $D$ \quad $\therefore M$
  \end{parts}

  \item Explain Mill's methods of agreement and difference with suitable examples.\pts{10}
  
  मिल की अन्वय तथा व्यतिरेक विधियों की उदाहरण सहित व्याख्या कीजिये।

  \item Explain any three syllogistic fallacies with examples.\pts{10}
  
  निरपेक्ष न्यायवाक्य के किन्हीं तीन दोषों की सोदाहरण व्याख्या कीजिये।

\end{enumerate}

\bigskip

\section*{SECTION -- B / खण्ड -- ब}
\noindent \textit{Note: Short Answer Type Questions. Attempt any \textbf{four} questions, each in 150 words. Each question carries 5 marks.}\pts{4 $\times$ 5 = 20}

\noindent \textit{किन्हीं चार प्रश्नों के उत्तर दीजिये, प्रत्येक उत्तर 150 शब्दों में हो। प्रत्येक प्रश्न 5 अंकों का है।}

\begin{enumerate}[label=\textbf{\arabic*.},leftmargin=*,start=6,itemsep=10pt]

  \item Distinguish between sentence and proposition.\pts{5}
  
  वाक्य एवं तर्कवाक्य में अन्तर स्पष्ट कीजिये।

  \item Discuss the distribution of terms in categorical propositions.\pts{5}
  
  निरपेक्ष तर्कवाक्यों में पदों की व्याप्ति की चर्चा कीजिये।

  \item Determine the validity/invalidity of ATI-1 according to Venn Diagram method.\pts{5}
  
  वेन आरेख विधि के अनुसार ATI-1 की वैधता/अवैधता निर्धारित कीजिये।

  \item What is argument by analogy? Discuss it with an example.\pts{5}
  
  सादृश्यानुमान युक्ति क्या है? उदाहरण सहित इसकी चर्चा कीजिये।

  \item Elucidate the traditional Square of Opposition with the help of a diagram. Discuss the logical relations that exist among the four standard categorical propositions.\pts{5}
  
  आरेख की सहायता से पारम्परिक विरोध-वर्ग को स्पष्ट कीजिये।

  \item Explain the principle of conversion in the light of Aristotelian categorical propositions.\pts{5}
  
  अरस्तू के निरपेक्ष तर्कवाक्यों के प्रकाश में परिवर्तन (Conversion) के नियम की व्याख्या कीजिये।

\end{enumerate}

\bigskip

\section*{SECTION -- C / खण्ड -- स}
\noindent \textit{Note: Very Short Answer Type Questions. Attempt all \textbf{ten} questions of this section in a maximum of two sentences each. Each question carries 2 marks.}\pts{10 $\times$ 2 = 20}

\noindent \textit{इस खण्ड के सभी दस प्रश्नों के उत्तर अधिकतम दो वाक्यों में दीजिये। प्रत्येक प्रश्न 2 अंकों का है।}

\begin{enumerate}[label=\textbf{\arabic*.},leftmargin=*,start=12,itemsep=8pt]
  \item 
  \begin{parts}
    \item What is the `Law of Excluded Middle'?\pts{2}
    
    `मध्यम-परिहार का नियम' क्या है?

    \item When E proposition is given false, what can we infer about A, I and O?\pts{2}
    
    जब E तर्कवाक्य असत्य दिया गया हो, तो A, I तथा O के विषय में क्या अनुमान लगाया जा सकता है?

    \item Write the symbolic form of: ``You will succeed if and only if you work hard.''\pts{2}
    
    प्रतीकात्मक रूप लिखिये: ``आप सफलता प्राप्त करेंगे यदि और केवल यदि आप कठिन परिश्रम करते हैं।''

    \item Is negation ($\sim$) a truth function?\pts{2}
    
    क्या निषेध ($\sim$) एक सत्यता फलन है?

    \item Name any two logical connectives and give their symbols.\pts{2}
    
    किन्हीं दो तार्किक संयोजकों के नाम तथा उनके प्रतीक दीजिये।

    \item Can we infer the truth value of I if A proposition is given false according to Traditional Square of Opposition?\pts{2}
    
    पारम्परिक विरोध-वर्ग के अनुसार यदि A तर्कवाक्य असत्य हो, तो क्या I का सत्यता मूल्य ज्ञात किया जा सकता है?

    \item Which term of the Categorical Syllogism connects the two premises?\pts{2}
    
    निरपेक्ष न्यायवाक्य का कौन-सा पद दो आधार वाक्यों को जोड़ता है?

    \item When is any argument invalid?\pts{2}
    
    कोई युक्ति कब अवैध होती है?

    \item Symbolise: ``Rohan and Mohan will not both be elected.'' ($R_x$: Rohan is elected, $M_x$: Mohan is elected)\pts{2}
    
    प्रतीकात्मक रूप: ``रोहन और मोहन दोनों निर्वाचित नहीं होंगे।''

    \item What is a tautologous statement?\pts{2}
    
    पुनरुक्ति (Tautology) कथन क्या होता है?
  \end{parts}
\end{enumerate}

\end{document}
